% Pursuing Stacks (cotes 134-2 to 134-8): archive cote 134-3, site page 57 (PDF page 58), typescript page 57.
% typescript. Draft compared: ps2.tex:2879-2946.
% AI-transcribed by Claude Opus 5.5 (claude-opus-5-5, adaptive thinking, effort high), Message Batches API, batch ps1 (msgbatch_01Rv5hCaNTtUt4JHx1yqoHPL), 2026-11-12; not yet checked by eye.
\scanpage{134-3/57}
\mspage{57}

To define a functor (*) , functorial in C for variable C, we only
have to choose a final element $e_a$ in each $\hand{\Delta}_a$ . For an element (a,p) of
$A_{/C}$ , we define
\[ \phi(a,p) = p(e_a) \quad , \]
with evident extension to arrows of $A_{/C}$\hand{\unsure{,}}(NB here it is important that the
$e_a$ be \uline{final} elements of the $\hand{\Delta}_a$ , initial elements for a change would'nt
do at all !), thus we get the functor $\phi$ , and functoriality with respect to
C is clear. Using the standart criterion of asphericity of a functor $\phi$ ,
we have to check that the categories
\[ (A_{/C})_{/x} \]
for $x \hand{\in}$ Ob C are aspheric, but one checks at once that the category above
is canonically isomorphic to $A_{/C'}$ , where $C' = C_{/x}$ . Thus asphericity
of $\phi$ for arbitrary C is equivalent to asphericity of $A_{/C}$ when C has a
final element, or equivalently, to asphericity of the element $\hand{\alpha}(C)$ in $\hand{A}\hat{}$
for such C. Now let $e_C$ be the final element of C, and consider the \xout{stan}
unique homotopy
\[ h : \hand{\Delta}_1 x C \hand{\to} C \]
between $id_C$ and the constant functor from C to C with value $e_C$ . Applying
the functor $\hand{\alpha}$ we get a corresponding homotopy in $A\hat{}$
\[ \hand{I} \times \hand{\alpha}(C) \hand{\to} \hand{\alpha}(C) \quad , \text{ where } I = \hand{\alpha}(\hand{\Delta}_1) \ , \]
between the identity map of $F = \hand{\alpha}(C)$ and a ``constant'' map $F \hand{\to} F$ . As we
assume that I is aspherical over the final element of $A\hat{}$, hence $IxF \hand{\to} F$
is a weak equivalence, it follows that $F \hand{\to} e$ is a weak equivalence, hen-%
ce F is aspheric as e is supposed to be aspheric, which was to be proved.

If we apply this proposition to $A=\hand{\Delta}$ and $n \hand{\to} \hand{\Delta}_n$ , we are still
reduced to check just T H 2) (which we did above looking at shuffles),
whereas T H 1) appears superfluous now - we get asphericity of the
elements $\hand{\Delta}_m x \hand{\Delta}_n$ in $A\hat{}$ as a consequence (but the homotopy argument used in \inserted{the}
prop. is essentially the same as the one used for proving T H 1 ). This
proposition applies equally to the case when $A = \hand{\widetilde{\Delta}}$ (category of non ordered
simplices $\hand{\widetilde{\Delta}}_n$ ), here the natural functor
\[ \hand{i} : \hand{\widetilde{\Delta} \longrightarrow} \text{(Cat)} \quad , \hand{a \mapsto \mathfrak{P}^*}(a) \]
is obtained by associating to any finite set a the combinatorial simplex
it defines, embodied by the ordered set $\hand{\mathfrak{P}^*}(a)$ of it's facettes of all di-%
mensions, which can be identified \hand{with} the ordered set of all non-empty sub-%
sets of the finite set a. As we know already $\hand{\widetilde{\Delta}}$ is a test category, all
that remains to be done is to check in this case that $I = \hand{\alpha}(\hand{\Delta}_1)$ is asphe-%
ric (hence aspheric over the final object of $A\hat{}$), Now this follows \xout{from}
\xout{the fact} again from a homotopy argument involving the unit segment substi-%
